\documentclass[11pt]{article}
\usepackage[a4paper,margin=18mm]{geometry}
\usepackage[no-math]{fontspec}
\setmainfont{Latin Modern Roman}
\usepackage{amsmath,amssymb,amsthm,xfrac,xspace,hyperref,graphicx,comment}
\excludecomment{DefTralics}

\providecommand{\bibliofont}{\small}
\newcommand{\ie}{i.e.,\xspace}
\newcommand{\eg}{e.g.,\xspace}
\newcommand{\etc}{etc.\xspace}
\newcommand{\cf}{cf.\xspace}
\newcommand{\resp}{resp.\xspace}
\newcommand{\smaller}{\small}
\newcommand{\Subsection}{\subsection}
\newcommand{\Subsubsection}{\subsubsection}
\newcommand{\skpt}{\smallskip}
\emergencystretch=3em
\numberwithin{equation}{section}\input{author-preamble.tex}
\begin{document}
\begin{center}{\Large Stable item 2054}\end{center}
\paragraph{Source and attribution.}
Jérémi Dardé, Julien Royer, \emph{Critical time for the observability of Kolmogorov-type equations}, Journal de l’École polytechnique — Mathématiques 8 (2021), publisher version, DOI 10.5802/jep.160. \href{https://jep.centre-mersenne.org/articles/10.5802/jep.160/}{Publisher article}. Authors' text: \href{https://creativecommons.org/licenses/by/4.0/}{CC BY 4.0}.
\paragraph{Editorial problem boundary.}
Prove only Proposition 3.6 below for the printed Dirichlet operator \(K_n=-\partial_{yy}+in q(y)^2\) on \(I=(-\ell_-,\ell_+)\), with \(q\in C^3(\overline I;\mathbf R)\), \(q(0)=0\) and strictly positive printed derivative bound. For each \(0<\gamma<q'(0)/\sqrt2\), prove the large-positive-integer uniform half-plane resolvent estimate and the exact zero-resolvent comparison to the harmonic model \(H_n\). The entire two-scale cut-off, harmonic/Airy parametrix and almost-orthogonal summation proof is retained. Observability times, Agmon boundary estimates and unrelated semigroup conclusions remain context only.
\paragraph{Difficulty (editorial): H2.}
Separate harmonic and Airy resolvent bounds do not control the variable-coefficient Dirichlet operator uniformly over the expanding spectral half-plane. The author builds two compatible localization scales, controls Taylor and commutator errors, and sums a growing family of local inverse operators using almost orthogonality. A further zero-resolvent comparison removes the cutoff and gives the required little-o model error.
\paragraph{Editorial transcription note.}
Original author language, mathematics and ordering are preserved. Publisher front matter is replaced by this independent article wrapper. Bibliography is recovered from matching cached publisher HTML, including its TeX formulas. No new proof or mathematical repair is supplied. Surrounding original results are retained for conservative context and are not extra selected corpus claims.
\clearpage
\input{author-body.tex}
\begin{thebibliography}{99}
\bibitem{Allonsius_Boyer_Morancey} D. Allonsius, F. Boyer \& M. Morancey - “Analysis of the null-controllability of degenerate parabolic systems of Grushin type via the moments method” , 2020, to appear in J. Evol. Equ.
\bibitem{Agmon85} S. Agmon - “Bounds on exponential decay of eigenfunctions of Schrödinger operators” , in Schrödinger operators (Como, 1984) , Lect. Notes in Math. , vol. 1159 , Springer, Berlin, 1985, p. 1-38
\bibitem{AmmarKhodja_Benabdallah_GonzalesBurgos_deTeresa} F. Ammar Khodja, A. Benabdallah, M. González-Burgos \& L. de Teresa - “New phenomena for the null controllability of parabolic systems: minimal time and geometrical dependence” , J. Math. Anal. Appl. 444 (2016) no. 2, p. 1071-1113
\bibitem{Bakri12} L. Bakri - “Quantitative uniqueness for Schrödinger operator” , Indiana Univ. Math. J. 61 (2012) no. 4, p. 1565-1580
\bibitem{Benabdallah_Boyer_Morancey} A. Benabdallah, F. Boyer \& M. Morancey - “A block moment method to handle spectral condensation phenomenon in parabolic control problems” , Ann. H. Lebesgue 3 (2020), p. 717-793
\bibitem{Beauchard_Cannarsa} K. Beauchard \& P. Cannarsa - “Heat equation on the Heisenberg group: observability and applications” , J. Differential Equations 262 (2017) no. 8, p. 4475-4521
\bibitem{Beauchard_Cannarsa_Guglielmi} K. Beauchard, P. Cannarsa \& R. Guglielmi - “Null controllability of Grushin-type operators in dimension two” , J. Eur. Math. Soc. (JEMS) 16 (2014) no. 1, p. 67–101
\bibitem{Beauchard_Darde_Ervedoza} K. Beauchard, J. Dardé \& S. Ervedoza - “Minimal time issues for the observability of Grushin-type equations” , Ann. Inst. Fourier (Grenoble) 70 (2020) no. 1, p. 247-312
\bibitem{Beauchard14} K. Beauchard - “Null controllability of Kolmogorov-type equations” , Math. Control Signals Systems 26 (2014) no. 1, p. 145-176
\bibitem{beauchard_henry} K. Beauchard, B. Helffer, R. Henry \& L. Robbiano - “Degenerate parabolic operators of Kolmogorov type with a geometric control condition” , ESAIM Control Optim. Calc. Var. 21 (2015) no. 2, p. 487-512
\bibitem{bardoslr92} C. Bardos, G. Lebeau \& J. Rauch - “Sharp sufficient conditions for the observation, control, and stabilization of waves from the boundary” , SIAM J. Control Optim. 30 (1992) no. 5, p. 1024-1065
\bibitem{Beauchard_Miller_Morancey} K. Beauchard, L. Miller \& M. Morancey - “2D Grushin-type equations: minimal time and null controllable data” , J. Differential Equations 259 (2015) no. 11, p. 5813-5845
\bibitem{Burq_Sun_19} N. Burq \& C. Sun - “Time optimal observability for Grushin Schrödinger equation” , 2019
\bibitem{Beauchard_Zuazua_09} K. Beauchard \& E. Zuazua - “Some controllability results for the 2D Kolmogorov equation” , Ann. Inst. H. Poincaré Anal. Non Linéaire 26 (2009) no. 5, p. 1793-1815
\bibitem{Cannarsa_Martinez_Vancostenoble} P. Cannarsa, P. Martinez \& J. Vancostenoble - Global Carleman estimates for degenerate parabolic operators with applications , Mem. Amer. Math. Soc. , vol. 239, no. 1133 , American Mathematical Society, Providence, RI, 2016
\bibitem{Duprez_Koenig} M. Duprez \& A. Koenig - “Control of the Grushin equation: non-rectangular control region and minimal time” , ESAIM Control Optim. Calc. Var. 26 (2020), article ID 3, 18 pages
\bibitem{Duprez} M. Duprez - “Controllability of a $2\times 2$ parabolic system by one force with space-dependent coupling term of order one” , ESAIM Control Optim. Calc. Var. 23 (2017) no. 4, p. 1473-1498
\bibitem{egorov} J. V. Egorov - “Some problems in the theory of optimal control” , Ž. Vyčisl. Mat. i Mat. Fiz. 3 (1963), p. 887-904
\bibitem{engel} K.-J. Engel \& R. Nagel - One-parameter semigroups for linear evolution equations , Graduate Texts in Math. , vol. 194 , Springer-Verlag, New York, 2000
\bibitem{Fursikov_Imanuvilov} A. V. Fursikov \& O. Y. Imanuvilov - Controllability of evolution equations , Lect. Notes Series , vol. 34 , Seoul National University, Seoul, 1996
\bibitem{fattorini_russel} H. O. Fattorini \& D. L. Russell - “Exact controllability theorems for linear parabolic equations in one space dimension” , Arch. Rational Mech. Anal. 43 (1971), p. 272-292
\bibitem{helffer1} B. Helffer - Semi-classical analysis for the Schrödinger operator and applications , Lect. Notes in Math. , vol. 1336 , Springer-Verlag, Berlin, 1988
\bibitem{Helffer11} B. Helffer - “On pseudo-spectral problems related to a time-dependent model in superconductivity with electric current” , Confluentes Math. 3 (2011) no. 2, p. 237-251
\bibitem{helffer} B. Helffer - Spectral theory and its applications , Cambridge Studies in Advanced Math. , vol. 139 , Cambridge University Press, Cambridge, 2013
\bibitem{henry} R. Henry - “On the semi-classical analysis of Schrödinger operators with purely imaginary electric potentials in a bounded domain” , 2014
\bibitem{HelfferSjo10} B. Helffer \& J. Sjöstrand - “From resolvent bounds to semi-group bounds” , 2010, Actes du colloque d’Évian 2009
\bibitem{HitrikSjoVio13} M. Hitrik, J. Sjöstrand \& J. Viola - “Resolvent estimates for elliptic quadratic differential operators” , Anal. PDE 6 (2013) no. 1, p. 181-196
\bibitem{kato} T. Kato - Perturbation theory for linear operators , Classics in Mathematics , Springer-Verlag, Berlin, 1980
\bibitem{Koenig17} A. Koenig - “Non-null-controllability of the Grushin operator in 2D” , Comptes Rendus Mathématique 355 (2017) no. 12, p. 1215-1235
\bibitem{koenig18} A. Koenig - “Lack of null-controllability for the fractional heat equation and related equations” , SIAM J. Control Optim. 58 (2020) no. 6, p. 3130-3160
\bibitem{KrejcirikRayRoySie17} D. Krejčiřík, N. Raymond, J. Royer \& P. Siegl - “Non-accretive Schrödinger operators and exponential decay of their eigenfunctions” , Israel J. Math. 221 (2017) no. 2, p. 779-802
\bibitem{KrejcirikSie15} D. Krejčiřík \& P. Siegl - “Elements of spectral theory without the spectral theorem” , in Non-selfadjoint operators in quantum physics , Wiley, Hoboken, NJ, 2015, p. 241-291
\bibitem{KrejcirikSieTatVio15} D. Krejčiřík, P. Siegl, M. Tater \& J. Viola - “Pseudospectra in non-Hermitian quantum mechanics” , J. Math. Phys. 56 (2015) no. 10, article ID 103513, 32 pages
\bibitem{Lissy19} P. Lissy - “A non-controllability result for the half-heat equation on the whole line based on the prolate spheroidal wave functions and its application to the Grushin equation” , 2019
\bibitem{LaurentLeautaud18} C. Laurent \& M. Léautaud - “Tunneling estimates and approximate controllability for hypoelliptic equations” , 2017, to appear in Mem. Amer. Math. Soc.
\bibitem{Lebeau_Robbiano} G. Lebeau \& L. Robbiano - “Contrôle exact de l’équation de la chaleur” , Comm. Partial Differential Equations 20 (1995) no. 1-2, p. 335-356
\bibitem{LetrouitSun20} C. Letrouit \& C. Sun - “Observability of Baouendi-Grushin-type equations through resolvent estimates” , 2020
\bibitem{raucht74} J. Rauch \& M. Taylor - “Exponential decay of solutions to hyperbolic equations in bounded domains” , Indiana Univ. Math. J. 24 (1974), p. 79-86
\bibitem{tucsnak} M. Tucsnak \& G. Weiss - Observation and control for operator semigroups , Birkhäuser Advanced Texts , Birkhäuser Verlag, Basel, 2009
\bibitem{ValleeSoares} O. Vallée \& M. Soares - Airy functions and applications to physics , Imperial College Press, London, 2010
\end{thebibliography}
\end{document}
